\documentclass[12pt,a4paper]{article}
\usepackage[utf8]{inputenc}
\usepackage[portuguese,english]{babel}
\usepackage{amsmath}
\usepackage{amssymb}
\usepackage{amsthm}
\usepackage{mathtools}
\usepackage{hyperref}

\newtheorem{theorem}{Theorem}
\newtheorem{corollary}{Corollary}[theorem]
\newtheorem{lemma}{Lemma}
\newtheorem{definition}{Definition}
\newtheorem{hypothesis}{Hypothesis}

\title{TCOP: Theorem of Primitive Operational Composition}
\author{}
\date{}

\begin{document}

\maketitle

\begin{abstract}
This document presents the formal statement and proof of the Theorem of Primitive Operational Composition (TCOP), which establishes that early non-phonetic graphic systems encode ontological classes and operational state transitions through finite graphical composition, implementing a closed algebra of validation rather than natural language.
\end{abstract}

\section{Introduction}

Early human graphic systems, including the Indus Valley script, Ancient Chinese characters, early Egyptian hieroglyphs, and the Phaistos Disk, share structural properties that distinguish them from phonetic writing systems. The TCOP formalizes these properties and demonstrates their mathematical equivalence to finite automata with attributes.

\section{Formal Statement}

\begin{theorem}[Primitive Operational Composition]
\label{thm:tcop}
Let $S = (B, O, \Sigma)$ be a primitive graphic system used for institutional, administrative, or ritual record-keeping, where:
\begin{itemize}
    \item $B$ is a finite set of base symbols (classes)
    \item $O$ is a finite set of graphical operators
    \item $\Sigma$ is the set of realized graphic states (compound symbols)
\end{itemize}

If the following conditions hold:
\begin{enumerate}
    \item (Recurrence) $B$ is finite and recurrent across the corpus
    \item (Systematicity) Base symbols admit systematic graphical modifications
    \item (Compositionality) Modifications are finite, recurrent, and compositional
    \item (Regularity) Symbol positions and co-occurrences exhibit statistical regularities
    \item (Non-phonetic dominance) No evidence of dominant phonetic encoding
\end{enumerate}

Then $S$ implements a finite operational algebra in which:
\begin{enumerate}
    \item Base symbols encode ontological classes
    \item Additional symbols encode state operators
    \item Meaning emerges from composition, not phonetic reading
    \item The system is formally equivalent to a finite automaton with attributes
\end{enumerate}
\end{theorem}

\section{Fundamental Hypotheses}

\begin{hypothesis}[Finiteness]
$|B| \ll |\Sigma|$, where $|B|$ is the number of base symbols and $|\Sigma|$ is the number of realized composite symbols.
\end{hypothesis}

\begin{hypothesis}[Composition]
Every $s \in \Sigma$ can be decomposed as:
\[
s = o_n \circ o_{n-1} \circ \cdots \circ o_1(b), \quad b \in B
\]
where $o_i \in O$ are operators applied sequentially to a base symbol $b$.
\end{hypothesis}

\begin{hypothesis}[Stability]
Operators do not create new classes, only new states within existing classes.
\end{hypothesis}

\begin{hypothesis}[Closure]
There exists a subset $O_T \subseteq O$ of terminal operators (closures) that mark completion states.
\end{hypothesis}

\begin{hypothesis}[Distribution]
Symbol frequencies follow a truncated Zipf-like hierarchy:
\[
f(r) \propto \frac{1}{r^\alpha}, \quad \alpha \in [0.8, 1.2]
\]
where $r$ is the rank and $f(r)$ is the frequency of the $r$-th most common symbol.
\end{hypothesis}

\section{Formal Proof}

\begin{proof}
We prove the theorem through a sequence of lemmas establishing the formal equivalence between primitive graphic systems satisfying the stated conditions and finite automata with attributes.

\begin{lemma}[Finite State Space]
Under Hypotheses 1 and 2, the system $S$ has a finite state space bounded by:
\[
|\Sigma| \leq |B| \cdot \prod_{i=1}^{k} (|O_i| + 1)
\]
where $k$ is the maximum depth of operator composition.
\end{lemma}

\begin{proof}[Proof of Lemma 1]
By Hypothesis 1 (Finiteness), $|B|$ is finite. By Hypothesis 2 (Composition), each symbol $s \in \Sigma$ is constructed by applying a finite sequence of operators to a base symbol. Since the system is used for record-keeping (not infinite generative language), the composition depth $k$ is bounded by practical constraints. Each position can have at most $|O_i| + 1$ choices (apply operator or skip). Therefore, the total number of possible states is bounded by the product.
\end{proof}

\begin{lemma}[Operational Closure]
Under Hypothesis 4, the system $S$ partitions into equivalence classes of reachable states, with terminal operators defining accepting states.
\end{lemma}

\begin{proof}[Proof of Lemma 2]
Define an equivalence relation $\sim$ on $\Sigma$ where $s_1 \sim s_2$ if they share the same base class $b \in B$. By Hypothesis 3 (Stability), operators preserve class membership. By Hypothesis 4 (Closure), terminal operators $O_T$ mark final states, creating a partition of $\Sigma$ into equivalence classes indexed by $B$ with defined terminal states.
\end{proof}

\begin{lemma}[Automaton Equivalence]
The system $S$ is equivalent to a deterministic finite automaton $M = (Q, \Sigma', \delta, q_0, F)$ with attribute functions, where:
\begin{itemize}
    \item $Q = B \times \mathcal{P}(O)$ (states are base classes with applied operators)
    \item $\Sigma'$ is the input alphabet (operators)
    \item $\delta: Q \times \Sigma' \to Q$ is the transition function
    \item $q_0 \in B$ (start with a base class)
    \item $F \subseteq Q$ (accepting states marked by terminal operators)
\end{itemize}
\end{lemma}

\begin{proof}[Proof of Lemma 3]
Construct the automaton as follows:
\begin{enumerate}
    \item States: Each state $(b, O')$ represents base class $b$ with applied operator set $O' \subseteq O$.
    \item Transitions: $\delta((b, O'), o) = (b, O' \cup \{o\})$ for $o \in O$ if $o$ is applicable.
    \item Initial state: Start with any $b \in B$ and $O' = \emptyset$.
    \item Accepting states: $F = \{(b, O') : O' \cap O_T \neq \emptyset\}$ (states with terminal operators).
\end{enumerate}

By Hypothesis 2, every symbol in $\Sigma$ corresponds to a unique state reachable through operator applications. By construction, this defines a finite automaton that accepts exactly the valid symbol sequences in $S$.
\end{proof}

\begin{lemma}[Non-linguistic Nature]
The system $S$ does not encode natural language in the linguistic sense.
\end{lemma}

\begin{proof}[Proof of Lemma 4]
Natural languages exhibit:
\begin{itemize}
    \item Unbounded recursion (context-free or context-sensitive grammars)
    \item Phonetic encoding with systematic sound-meaning correspondences
    \item Narrative structures with arbitrary length
    \item Zipf distribution across thousands of distinct lexical items
\end{itemize}

System $S$ exhibits:
\begin{itemize}
    \item Finite state space (Lemma 1)
    \item No evidence of phonetic dominance (Theorem condition 5)
    \item Short inscriptions (practical record-keeping)
    \item Truncated distribution over hundreds of states (Hypothesis 5)
\end{itemize}

The structural properties of $S$ are incompatible with natural language encoding. Instead, $S$ implements a finite operational algebra for state validation and record-keeping.
\end{proof}

Combining Lemmas 1-4, we have established that system $S$ satisfying the theorem's conditions:
\begin{enumerate}
    \item Has finite state space (Lemma 1)
    \item Partitions into operational classes (Lemma 2)
    \item Is equivalent to a finite automaton with attributes (Lemma 3)
    \item Does not encode natural language (Lemma 4)
\end{enumerate}

Therefore, $S$ implements a closed algebra of validation through finite graphical composition.
\end{proof}

\section{Corollaries}

\begin{corollary}[Ancient Chinese Application]
In Ancient Chinese, the radical 人 (rén) does not possess fixed semantic meaning but defines an agent class whose state varies according to added graphical operators (duplication, force, authority, internalization, etc.).

Formally:
\[
\text{State} = \text{Human} \circ \sum \text{Operators}
\]
\end{corollary}

\begin{corollary}[Harappan Application]
Harappan symbols are not words but validated states resulting from applying finite operators to base classes. This explains:
\begin{itemize}
    \item Short inscriptions (typically 5-10 symbols)
    \item Absence of narrative structure
    \item Recurrence of terminal signs
    \item System saturation at approximately 417 distinct states
\end{itemize}
\end{corollary}

\begin{corollary}[Cognitive Universal]
Systems such as Harappan, early Egyptian hieroglyphs, Ancient Chinese, and the Phaistos Disk share a common cognitive substrate in which graphical operators precede and survive phonetic languages.
\end{corollary}

\begin{corollary}[Mathematical Limit]
The mathematics implemented by such systems is necessarily finite, relational, and state-based, bounded by:
\[
N_{\text{finite}} + \text{relations} + \text{partial orders}
\]

Symbolic algebra and infinite mathematics require a different historical paradigm.
\end{corollary}

\section{Modern Interpretation}

The TCOP implies that primitive graphic systems are formally equivalent to:
\begin{itemize}
    \item Workflows
    \item Ledgers
    \item Finite automata
    \item JSON records
    \item Validation protocols
\end{itemize}

They are \textbf{not} equivalent to literary texts or phonetic writing systems.

\section{Limitations}

The TCOP does \textbf{not} claim:
\begin{itemize}
    \item Direct common historical origin
    \item Linguistic identity across systems
    \item Continuous cultural transmission
    \item Universality without exceptions
\end{itemize}

It claims only:
\begin{itemize}
    \item Structural equivalence under common cognitive and functional constraints
\end{itemize}

\section{Defense Against Classical Criticisms}

\subsection{Linguistic Criticism}

\textbf{Criticism:} "All writing systems encode language, therefore these systems must encode language."

\textbf{Response:} This assumes the conclusion. The TCOP demonstrates that graphic systems with the specified properties implement operational algebras, not phonetic encoding. The absence of:
\begin{itemize}
    \item Long narrative texts
    \item Systematic phonetic correspondences
    \item Recursive grammatical structures
\end{itemize}
supports the operational interpretation over the linguistic interpretation.

\subsection{Archaeological Criticism}

\textbf{Criticism:} "Without bilingual texts, we cannot know what these systems encode."

\textbf{Response:} The TCOP is based on \textit{structural analysis}, not translation. The mathematical properties (finite state space, compositional operators, statistical regularities) can be determined without decipherment. The theorem predicts specific testable properties that distinguish operational systems from linguistic systems.

\subsection{Historical Criticism}

\textbf{Criticism:} "These systems evolved into phonetic writing, so they must have been proto-phonetic."

\textbf{Response:} Evolution does not imply identity. The TCOP describes a specific class of systems at a specific historical moment. That some systems later developed phonetic components does not invalidate the operational nature of their earlier forms. The theorem explicitly allows for historical transitions.

\subsection{Universalist Criticism}

\textbf{Criticism:} "The theorem claims universal cognitive constraints without evidence."

\textbf{Response:} The TCOP claims \textit{structural equivalence under specific constraints}, not universal cognitive laws. The constraints (finite memory, institutional record-keeping, visual distinctiveness) are empirically grounded. The theorem makes testable predictions about any system satisfying these constraints.

\section{Conclusion}

The Theorem of Primitive Operational Composition provides a mathematical framework for understanding early graphic systems as operational algebras rather than phonetic encodings. This framework:
\begin{itemize}
    \item Explains structural properties previously considered anomalous
    \item Makes testable predictions about undeciphered scripts
    \item Provides formal tools for comparative analysis
    \item Establishes bounds on the mathematical sophistication of such systems
\end{itemize}

\end{document}
